% Options for packages loaded elsewhere
\PassOptionsToPackage{unicode}{hyperref}
\PassOptionsToPackage{hyphens}{url}
\documentclass[
  ignorenonframetext,
]{beamer}
\newif\ifbibliography
\usepackage{pgfpages}
\setbeamertemplate{caption}[numbered]
\setbeamertemplate{caption label separator}{: }
\setbeamercolor{caption name}{fg=normal text.fg}
\beamertemplatenavigationsymbolsempty
% remove section numbering
\setbeamertemplate{part page}{
  \centering
  \begin{beamercolorbox}[sep=16pt,center]{part title}
    \usebeamerfont{part title}\insertpart\par
  \end{beamercolorbox}
}
\setbeamertemplate{section page}{
  \centering
  \begin{beamercolorbox}[sep=12pt,center]{section title}
    \usebeamerfont{section title}\insertsection\par
  \end{beamercolorbox}
}
\setbeamertemplate{subsection page}{
  \centering
  \begin{beamercolorbox}[sep=8pt,center]{subsection title}
    \usebeamerfont{subsection title}\insertsubsection\par
  \end{beamercolorbox}
}
% Prevent slide breaks in the middle of a paragraph
\widowpenalties 1 10000
\raggedbottom
\AtBeginPart{
  \frame{\partpage}
}
\AtBeginSection{
  \ifbibliography
  \else
    \frame{\sectionpage}
  \fi
}
\AtBeginSubsection{
  \frame{\subsectionpage}
}
\usepackage{iftex}
\ifPDFTeX
  \usepackage[T1]{fontenc}
  \usepackage[utf8]{inputenc}
  \usepackage{textcomp} % provide euro and other symbols
\else % if luatex or xetex
  \usepackage{unicode-math} % this also loads fontspec
  \defaultfontfeatures{Scale=MatchLowercase}
  \defaultfontfeatures[\rmfamily]{Ligatures=TeX,Scale=1}
\fi
\usepackage{lmodern}
\ifPDFTeX\else
  % xetex/luatex font selection
\fi
% Use upquote if available, for straight quotes in verbatim environments
\IfFileExists{upquote.sty}{\usepackage{upquote}}{}
\IfFileExists{microtype.sty}{% use microtype if available
  \usepackage[]{microtype}
  \UseMicrotypeSet[protrusion]{basicmath} % disable protrusion for tt fonts
}{}
\makeatletter
\@ifundefined{KOMAClassName}{% if non-KOMA class
  \IfFileExists{parskip.sty}{%
    \usepackage{parskip}
  }{% else
    \setlength{\parindent}{0pt}
    \setlength{\parskip}{6pt plus 2pt minus 1pt}}
}{% if KOMA class
  \KOMAoptions{parskip=half}}
\makeatother
\usepackage{longtable,booktabs,array}
\usepackage{caption}
\captionsetup[table]{skip=6pt}
\newcounter{none} % for unnumbered tables
\usepackage{calc} % for calculating minipage widths
\usepackage{caption}
% Make caption package work with longtable
\makeatletter
\def\fnum@table{\tablename~\thetable}
\makeatother
\setlength{\emergencystretch}{3em} % prevent overfull lines
\providecommand{\tightlist}{%
  \setlength{\itemsep}{0pt}\setlength{\parskip}{0pt}}
% Manuscript macro/package fallbacks (newcommand -> providecommand).
% Core mathematical packages
\usepackage{amsmath,amssymb,amsthm}
\usepackage{mathtools}
% Keep long deployment prose and references within the text block.
% Algorithm formatting
% Tables
\usepackage{booktabs}
\usepackage{multirow}
% Graphics and floats
% Typography and layout
% Page layout (slightly smaller margins for academic journal style)
% Running headers/footers
% Caption formatting
% Colored boxes for callouts
% Cross-referencing with smart naming
% Hyperlink styling
% Configure cleveref for custom environments
% Theorem environments with consistent numbering
% Code listings and raw implementation examples in the manuscript
\usepackage{listings}
\lstset{basicstyle=\ttfamily\small,breaklines=true,columns=fullflexible}
% Math operators
\DeclareMathOperator*{\argmax}{arg\,max}
\DeclareMathOperator*{\argmin}{arg\,min}
\DeclareMathOperator{\sign}{sign}
\DeclareMathOperator{\KL}{KL}
\DeclareMathOperator{\Tr}{Tr}
% Custom commands for notation consistency
\providecommand{\calA}{\mathcal{A}}
\providecommand{\calB}{\mathcal{B}}
\providecommand{\calC}{\mathcal{C}}
\providecommand{\calD}{\mathcal{D}}
\providecommand{\calF}{\mathcal{F}}
\providecommand{\calG}{\mathcal{G}}
\providecommand{\calH}{\mathcal{H}}
\providecommand{\calI}{\mathcal{I}}
\providecommand{\calM}{\mathcal{M}}
\providecommand{\calO}{\mathcal{O}}
\providecommand{\calP}{\mathcal{P}}
\providecommand{\calS}{\mathcal{S}}
\providecommand{\calT}{\mathcal{T}}
\providecommand{\calW}{\mathcal{W}}
% Note: \Phi is a standard LaTeX command, do not redefine
% Trust notation
\providecommand{\trust}[2]{\mathcal{T}_{#1 \to #2}}
\providecommand{\trustt}[3]{\mathcal{T}_{#1 \to #2}^{#3}}
% Belief notation
\providecommand{\belief}[2]{\mathcal{B}_{#1}(#2)}
\providecommand{\belieft}[3]{\mathcal{B}_{#1}^{#2}(#3)}
% Cognitive state
\providecommand{\cogstate}[1]{\sigma_{#1}}
% Attack notation
\providecommand{\attack}[1]{\mathcal{A}_{#1}}
\providecommand{\adversary}[1]{\Omega_{#1}}
% Defense notation
\providecommand{\firewall}{\mathcal{F}}
\providecommand{\sandbox}{\mathcal{S}_{box}}
% Probability and expectation
\providecommand{\E}{\mathbb{E}}
\providecommand{\Prob}{\mathbb{P}}
\providecommand{\indicator}{\mathbb{1}}
% QED symbol
\providecommand{\qedsymbol}{$\blacksquare$}
% Front matter opens on the cover/title page (no list of figures/tables precedes it).

% Slide layout defaults for warning-clean scientific decks.
\usepackage{etoolbox}
\IfFileExists{xurl.sty}{\usepackage{xurl}}{}
\IfFileExists{seqsplit.sty}{\usepackage{seqsplit}}{\newcommand{\seqsplit}[1]{#1}}
\protected\def\breakseq#1{\seqsplit{#1}}
\protected\def\breaktt#1{\begingroup\ttfamily\seqsplit{#1}\endgroup}
\setlength{\emergencystretch}{6em}
\tolerance=5000
\hbadness=10000
\hfuzz=1pt
\setlength{\tabcolsep}{2pt}
\AtBeginEnvironment{longtable}{\tiny\renewcommand{\arraystretch}{0.86}\setlength{\tabcolsep}{1pt}}
\AtBeginEnvironment{tabular}{\tiny\renewcommand{\arraystretch}{0.86}\setlength{\tabcolsep}{1pt}}
\AtBeginEnvironment{equation}{\tiny}
\AtBeginEnvironment{equation*}{\tiny}
\AtBeginEnvironment{align}{\tiny}
\AtBeginEnvironment{align*}{\tiny}
\AtBeginEnvironment{itemize}{\footnotesize}
\AtBeginEnvironment{enumerate}{\footnotesize}
\AtBeginEnvironment{description}{\footnotesize}
\setbeamerfont{caption}{size=\tiny}
\setbeamerfont{caption name}{size=\tiny}
\setbeamerfont{normal text}{size=\small}
\setbeamerfont{frametitle}{size=\small}
\setbeamerfont{section title}{size=\footnotesize}
\setbeamerfont{subsection title}{size=\footnotesize}
\setbeamertemplate{section page}{%
  \centering
  \begin{beamercolorbox}[sep=12pt,center,wd=\paperwidth]{section title}
    \parbox{0.86\paperwidth}{\centering\usebeamerfont{section title}\insertsection\par}
  \end{beamercolorbox}
}
\setbeamertemplate{subsection page}{%
  \centering
  \begin{beamercolorbox}[sep=8pt,center,wd=\paperwidth]{subsection title}
    \parbox{0.86\paperwidth}{\centering\usebeamerfont{subsection title}\insertsubsection\par}
  \end{beamercolorbox}
}
\setlength{\abovecaptionskip}{2pt}
\setlength{\belowcaptionskip}{0pt}

% Natbib and cross-reference fallbacks — slides don't load natbib
% or cleveref, but combined-PDF manuscript prose may emit these
% commands. The fallback renders citations as a bracketed key list
% and cross-references as detokenized labels so slides don't fail on
% undefined control sequences or raw underscores. \providecommand is
% a no-op if packages load later.
\providecommand{\citep}[1]{[#1]}
\providecommand{\citet}[1]{#1}
\providecommand{\citealp}[1]{#1}
\providecommand{\citeauthor}[1]{#1}
\providecommand{\citeyear}[1]{#1}
\providecommand{\cref}[1]{\texttt{\detokenize{#1}}}
\providecommand{\Cref}[1]{\texttt{\detokenize{#1}}}
\providecommand{\crefrange}[2]{\texttt{\detokenize{#1}}--\texttt{\detokenize{#2}}}
\providecommand{\Crefrange}[2]{\texttt{\detokenize{#1}}--\texttt{\detokenize{#2}}}
\providecommand{\paragraph}[1]{\textbf{#1}\ }

\newtheorem{proposition}{Proposition}
\newtheorem{hypothesis}{Hypothesis}
\newtheorem{remark}[theorem]{Remark}
\newtheorem{axiom}{Axiom}
\newtheorem{property}{Property}
\usepackage{bookmark}
\IfFileExists{xurl.sty}{\usepackage{xurl}}{} % add URL line breaks if available
\urlstyle{same}
% fallback for those not using the hyperref driver hyperxmp:
\makeatletter
\@ifundefined{xmpquote}{\newcommand{\xmpquote}[1]{#1}}{}
\makeatother
\hypersetup{
  hidelinks,
  pdfcreator={LaTeX via pandoc}}

\author{\texorpdfstring{}{}}
\date{}

\begin{document}

\begin{frame}
\newpage
\end{frame}

\section{Notation Reference}\label{sec:notation-reference}

\begin{frame}[allowframebreaks]{Notation Reference}
This paper intentionally minimizes mathematical notation to maximize
accessibility. Where notation is used, it follows the Cognitive
Integrity Framework (CIF) formal specification defined in Part 1 of this
series.
\end{frame}

\subsection{Minimal Notation Used}\label{minimal-notation-used}

\begin{frame}[allowframebreaks]{Minimal Notation Used}
{\def\LTcaptype{none} % do not increment counter
\begin{longtable}[]{@{}
  >{\raggedright\arraybackslash}p{(\linewidth - 4\tabcolsep) * \real{0.2424}}
  >{\raggedright\arraybackslash}p{(\linewidth - 4\tabcolsep) * \real{0.2727}}
  >{\raggedright\arraybackslash}p{(\linewidth - 4\tabcolsep) * \real{0.4848}}@{}}
\toprule\noalign{}
\begin{minipage}[b]{\linewidth}\raggedright
Symbol
\end{minipage} & \begin{minipage}[b]{\linewidth}\raggedright
Meaning
\end{minipage} & \begin{minipage}[b]{\linewidth}\raggedright
Plain Language
\end{minipage} \\
\midrule\noalign{}
\endhead
δ & Trust decay factor & ``Delegated trust decreases by this factor at
each step'' \\
n & Agent count & ``Number of agents in the system'' \\
f & Byzantine agents & ``Maximum number of malicious agents
tolerated'' \\
\([A]\) & Design Matrix & Maps Functional Requirements to Defense
Provisions \\
\(\{FR\}\) & Functional Requirements & What the system must protect \\
\(\{DP\}\) & Defense Provisions & What CIF mechanisms provide \\
\(\text{INV}_k\) & Individual invariant predicate & ``A hard rule the
system checks at runtime; Part 1 writes this \(I_k\)'' \\
\(\Omega_k\) & Adversary class \(k\) & Capability tier: 1=external
(input control), 2=peripheral (tool/data channels), 3=agent-level
(single compromised agent), 4=coordination (inter-agent channels),
5=systemic (orchestrator) \\
\bottomrule\noalign{}
\end{longtable}
}
\end{frame}

\subsection{CIF-AD-OODA Notation}\label{cif-ad-ooda-notation}

\begin{frame}[allowframebreaks]{CIF-AD-OODA Notation}
The cross-domain analysis (Sections 9c--9l) uses the CIF-AD-OODA
methodology:

\end{frame}

\begin{frame}[allowframebreaks]{CIF-AD-OODA Notation}
\begin{itemize}
\tightlist
\item
  \textbf{Design Matrix} \([A]\): A matrix where rows represent
  Functional Requirements (\(FR\)) and columns represent Defense
  Provisions (\(DP\)). Each entry \(A_{ij}\) indicates whether defense
  \(j\) covers requirement \(i\).
\item
  \textbf{Transient Coupling} \([A']\): The coupling matrix during an
  active attack, showing which defenses are bypassed.
\item
  \textbf{Adversary Classes} \(\Omega_1\)--\(\Omega_5\): Five adversary
  classes from external input control (\(\Omega_1\)) through peripheral
  tool/data-channel compromise (\(\Omega_2\)), single compromised agents
  (\(\Omega_3\)), coordination-channel attacks (\(\Omega_4\)), and
  systemic orchestrator compromise (\(\Omega_5\)), as defined in Part
  1's threat model.
\end{itemize}
\end{frame}

\begin{frame}[allowframebreaks]{Three Universal Attack Patterns}
\protect\phantomsection\label{three-universal-attack-patterns}
Across all ten domains, attacks reduce to three canonical patterns:

\end{frame}

\begin{frame}[allowframebreaks]{Three Universal Attack Patterns}
{\def\LTcaptype{none} % do not increment counter
\begin{longtable}[]{@{}
  >{\raggedright\arraybackslash}p{(\linewidth - 4\tabcolsep) * \real{0.2903}}
  >{\raggedright\arraybackslash}p{(\linewidth - 4\tabcolsep) * \real{0.4194}}
  >{\raggedright\arraybackslash}p{(\linewidth - 4\tabcolsep) * \real{0.2903}}@{}}
\toprule\noalign{}
\begin{minipage}[b]{\linewidth}\raggedright
Pattern
\end{minipage} & \begin{minipage}[b]{\linewidth}\raggedright
Description
\end{minipage} & \begin{minipage}[b]{\linewidth}\raggedright
Defense
\end{minipage} \\
\midrule\noalign{}
\endhead
FR Polarity Inversion & Attacker flips a Functional Requirement's sign
(e.g., ``don't share secrets'' → ``share secrets'') & Cognitive Firewall
+ Belief Sandbox \\
Constraint Relaxation & Attacker weakens a safety constraint's boundary
& Invariant Monitor + Tripwire \\
Context Boundary Violation & Attacker exploits scope leakage between
agent contexts & Provenance Tracking + Trust Calculus \\
\bottomrule\noalign{}
\end{longtable}
}
\end{frame}

\subsection{Trust Decay Explanation}\label{trust-decay-explanation}

\begin{frame}[allowframebreaks]{Trust Decay Explanation}
The symbol δ is a parameter, not a universal constant. For an
illustrative example, δ = 0.8 means:

\end{frame}

\begin{frame}[allowframebreaks]{Trust Decay Explanation}
\begin{itemize}
\tightlist
\item
  Direct trust: 100\% of assigned value
\item
  One delegation: 80\% of source trust
\item
  Two delegations: 64\% of source trust
\item
  Three delegations: 51.2\% of source trust
\end{itemize}

\end{frame}

\begin{frame}[allowframebreaks]{Trust Decay Explanation}
The executable Part 3 deployment profiles use δ = 0.80 for balanced
operation and δ = 0.60 for high assurance. Those are implementation
defaults, not a claim that either value is optimal for every threat
model. A lower δ means faster decay, providing more security against
long delegation chains while limiting delegation utility.
\end{frame}

\subsection{Byzantine Tolerance
Explanation}\label{byzantine-tolerance-explanation}

\begin{frame}[allowframebreaks]{Byzantine Tolerance Explanation}
When we say n \ensuremath{\ge} 3f + 1:

\end{frame}

\begin{frame}[allowframebreaks]{Byzantine Tolerance Explanation}
\begin{itemize}
\tightlist
\item
  To tolerate 1 malicious agent, need at least 4 agents
\item
  To tolerate 2 malicious agents, need at least 7 agents
\item
  To tolerate 3 malicious agents, need at least 10 agents
\end{itemize}
\end{frame}

\subsection{Full Notation Reference}\label{full-notation-reference}

\begin{frame}[fragile,allowframebreaks]{Full Notation Reference}
For complete formal definitions of all CIF notation, see Part 1
supplementary \textbf{S03}
(\breaktt{cogsec\_multiagent\_1\_theory/manuscript/S03\_notation.md} in
this repository's cognitive\_integrity program tree).

The Part 1 specification uses on the order of 100 symbols covering:

\end{frame}

\begin{frame}[fragile,allowframebreaks]{Full Notation Reference}
\begin{itemize}
\tightlist
\item
  Agent cognitive state
\item
  Trust calculus operations
\item
  Defense mechanism parameters
\item
  Consensus and coordination
\item
  Information-theoretic bounds
\end{itemize}
\end{frame}

\end{document}
